% ECONOMICS_PROBLEM: {"id": "R31-548", "title": "Housing-supply responses to reform", "jel_code": "R31", "source_id": "OP-0460"}
% !TeX program = xelatex
\documentclass[11pt,a4paper]{article}
\usepackage[margin=23mm]{geometry}
\usepackage{fontspec}
\setmainfont{Latin Modern Roman}
\usepackage{amsmath,amssymb,mathtools}
\usepackage{xcolor,enumitem,booktabs,longtable,array,xurl,hyperref}
\definecolor{muted}{HTML}{53636C}
\hypersetup{colorlinks=true,urlcolor=blue,linkcolor=blue}
\setlength{\parindent}{0pt}
\setlength{\parskip}{5pt}
\newcommand{\R}{\mathbb R}
\newcommand{\E}{\mathbb E}
\newcommand{\N}{\mathbb N}
\newcommand{\ind}{\mathbf 1}
\newcommand{\Var}{\operatorname{Var}}
\newcommand{\Cov}{\operatorname{Cov}}
\newcommand{\supp}{\operatorname{supp}}
\DeclareMathOperator*{\argmin}{arg\,min}
\DeclareMathOperator*{\argmax}{arg\,max}
\newcommand{\entrymeta}[3]{{\sffamily\small\color{muted}\textbf{\texttt{#1}}\quad|\quad #2\par #3\par}}
\newcommand{\Problemref}[1]{\texttt{#1}}
\newcommand{\JELref}[2]{\texttt{#2} (JEL \texttt{#1})}
\begin{document}
\section*{Housing-supply responses to reform}
\textbf{Problem ID:} \texttt{R31-548}\quad \textbf{JEL:} \texttt{R31}\par
% BEGIN ECONOMICS STATEMENT
\label{problem:OP-0460}
{\sffamily\small\color{muted}Original subject group: Urban, housing and spatial economics\par}
\label{definitions:problem:30007742}
\textbf{Audit classification:} Published research agenda. Checked source identity and appropriate agenda/background support; expanded intervention, units, population, definedness and causal restrictions. Exact targets and variants are editorial.
\subsection*{Definitions and precise research target}
\textit{Editorial formalization.} Consider metropolitan areas in one country adopting a specified reform that increases permitted residential density. Let \(a=1\) denote implementation and \(a=0\) continuation of prior rules. Define \(H_{mt}(a)\) as habitable housing units, \(R_{mt}(a)\) quality-adjusted real rent, and \(t=0\) as the reform year in metropolitan area \(m\). For every \(h\in\{1,\ldots,15\}\), define
\[\tau_H(h)=E[\log H_{m,h}(1)-\log H_{m,h}(0)],\qquad\tau_R(h)=E[\log R_{m,h}(1)-\log R_{m,h}(0)].\]
Require positive housing stocks and rents under both policies and integrable log contrasts. Identify these dynamic responses using randomized rollout or a justified comparison design, allowing construction delays, demolition, redevelopment, and infrastructure constraints. Record actual permitting and enforcement, since statutory eligibility does not equal realized construction. Keep baseline metropolitan boundaries and report spillovers to adjoining areas. An estimate at one horizon does not identify a complete supply elasticity: separately specify demand shocks and any price-based elasticity target. Determine which local conditions alter response speed and transport estimates only to areas with adequate covariate overlap.

Let \(t_m\) be the announcement or implementation date chosen before analysis, and define event time \(h=t-t_m\). \(H_{mh}(a)\) is the number of habitable units inside fixed baseline boundaries, including vacancy under a declared convention; \(R_{mh}(a)\) is an independently quality-adjusted real rental-service price. Both are positive and log contrasts integrable. Reform is the complete zoning/permitting/enforcement package; ancillary infrastructure is held fixed or included in the treatment. A dynamic supply elasticity instead concerns \(\partial\log H_{mh}/\partial\log R\) under a defined exogenous price/demand shift, and is not \(\tau_H/\tau_R\) without further restrictions. Interference across metros, endogenous adoption and differential construction trends require a design or model. Policy effect profiles and market elasticities are distinct targets. All expectations use a stated baseline population and probability law. Identification means every admissible data-generating model with the same observed-data law yields the same displayed target; otherwise report the identified set. Exact equations, horizons and assumptions are editorial, rather than source-given canonical conjectures.
\subsection*{Variants and assumption changes}
\begin{enumerate}
\item \textbf{Editorial redevelopment variant.} Decompose stock changes into greenfield building, redevelopment, demolition and conversion under one stock-flow identity. Permits are not completed units.
\item \textbf{Editorial quality variant.} Replace unit counts by independently adjusted service flows and compare size/maintenance effects. Unit growth need not imply proportional service growth.
\end{enumerate}
\subsection*{References and source audit}
\textbf{1.} Explicit agenda on dynamic supply, construction productivity, and richer migration models; exact targets are editorial. Locator: Handbook of Regional and Urban Economics, Volume 6, Chapter 6; Section 7, pp.448–449. Full text or primary publication page inspected. URL: \url{https://realestate.wharton.upenn.edu/wp-content/uploads/2025/12/BaumSnow_Duranton_bc_2025-003.pdf}.
\begin{enumerate}\raggedright
\item\hypertarget{definitions:source:30007742:1}{} Nathaniel Baum-Snow; Gilles Duranton. 2025. \emph{Housing supply and housing affordability}. Handbook of Regional and Urban Economics, Volume 6, Chapter 6; Section 7, pp.448–449.
\url{https://realestate.wharton.upenn.edu/wp-content/uploads/2025/12/BaumSnow_Duranton_bc_2025-003.pdf}.
DOI: \url{https://doi.org/10.1016/bs.hesreg.2025.05.003}.
\end{enumerate}

\subsection*{Common conventions: Source mathematical conventions}

These conventions apply to each entry unless explicitly replaced.

\textbf{Models and identification.} A primitive model \(m\) specifies all probability laws, feasible choices, information, equilibrium and measurement rules used by its target. The admissible class \(\mathcal M\) is fixed independently of outcomes. If \(P_O(m)\) is its observable probability law and \(\tau(m)\) the target, its identified set at observable law \(P\) is
\[
 \mathcal I_\tau(P)=\{\tau(m):m\in\mathcal M,\ P_O(m)=P\}.
\]
Point identification means that this set is a singleton. An empty set signals incompatibility of data and assumptions. Coordinate bounds need not describe the joint set. Bounds are sharp only if every retained target is attainable or arbitrarily approachable by compatible models. Finite bounds require explicit restrictions on unobserved quantities; unrestricted confounding, hidden wealth, extrapolation or arbitrary equilibrium selection can yield uninformative sets.

\textbf{Causal interventions.} A potential outcome \(Y_i(p)\) is the outcome under a complete policy regime \(p\), including timing, financing, information and market spillovers. The target population and units are fixed before treatment. With interference, use the complete assignment vector or a justified exposure mapping; individual no-interference notation alone cannot identify an economy-wide intervention. A difference of mean potential outcomes does not require the cross-world joint distribution. Conditioning on post-treatment migration, survival, enrollment or employment changes the estimand and needs separate justification.

\textbf{Statistical statements.} An estimator, confidence set, forecast or test specifies a sampling law, model class and loss/coverage criterion. Uniform \((1-\alpha)\)-coverage means \(\inf_{m\in\mathcal M}\Pr_m\{\tau(m)\in C_n\}\geq1-\alpha\), or an explicitly asymptotic version. The whole parameter space trivially has coverage, so informativeness requires a stated width, risk or power objective. A forecasting improvement is relative to a named benchmark, information set and evaluation population.

\textbf{Equilibrium and optimization.} An equilibrium comprises feasible plans, optimality, beliefs where needed, budget constraints and all market/resource clearing. The primitives or a stated benchmark define each correspondence \(\mathcal E(p;m)\). Nonemptiness is assumed or proved, never inferred just from compact policy sets. A selected-equilibrium target uses a declared selection rule; a robust target quantifies over all permitted equilibria. Use suprema or infima unless attainment is established. An \(\varepsilon\)-optimal policy is feasible and within \(\varepsilon\) of the finite optimum value. Report an empty feasible set separately.

\textbf{Welfare and accounting.} Normative utility, interpersonal normalization, social weights, numeraire, discounting and the use of public receipts are primitives. Behavioral utility need not coincide with the welfare criterion. Household welfare includes ownership income and rents once; adding profits again to compensating variation generally double-counts them. A monetary equivalent is defined by an explicit transfer and price convention, with monotonicity and range conditions. Average effects, mechanism contributions, transfers and welfare losses are different targets. Infinite sums require convergence and dynamic feasible plans require terminal or no-Ponzi conditions.

\textbf{Source versions.} A working-paper date, revision date and journal publication year are distinguished. A DOI for a journal version is not automatically the DOI of an earlier manuscript. Locators refer to the stated version and printed pages where specified. A metadata-only check does not verify a claim about a full-text conclusion. Every entry's source subsection narrows the provenance claim accordingly.
% END ECONOMICS STATEMENT
\end{document}
